\paper{非数学类}

\begin{problems}
\item[一、]（本题满分 24 分，共 4 小题，每小题 6 分）填空题：
\begin{enumerate}
  \item 设 $\alpha\in(0,1)$，则 $\displaystyle\lim_{n\to+\infty}\left[(n+1)^{\alpha}-n^{\alpha}\right]=\underline{\qquad\qquad}$．
  \item 若曲线 $y=y(x)$ 由 $\begin{cases}x=t+\cos t,\\ \ee^{y}+ty+\sin t=1\end{cases}$ 确定，则此曲线在 $t=0$ 对应点处的切线方程为 $\underline{\qquad\qquad}$．
  \item $\displaystyle\int\frac{\ln\left(x+\sqrt{1+x^{2}}\right)}{\left(1+x^{2}\right)^{3/2}}\dd x=\underline{\qquad\qquad}$．
  \item $\displaystyle\lim_{x\to 0}\frac{1-\cos x\sqrt{\cos 2x}\sqrt[3]{\cos 3x}}{x^{2}}=\underline{\qquad\qquad}$．
\end{enumerate}

\item[二、]（本题满分 8 分）设函数 $f(t)$ 在 $t\ne 0$ 时一阶连续可导，且 $f(1)=0$，求函数 $f\left(x^{2}-y^{2}\right)$，使得曲线积分
\[
\int_{L}y\left[2-f\left(x^{2}-y^{2}\right)\right]\dd x+xf\left(x^{2}-y^{2}\right)\dd y
\]
与路径无关，其中 $L$ 为任一不与直线 $y=\pm x$ 相交的分段光滑曲线．

\item[三、]（本题满分 14 分）设 $f(x)$ 在区间 $[0,1]$ 上连续，且 $1\le f(x)\le 3$．证明：
\[
1\le\int_{0}^{1}f(x)\dd x\int_{0}^{1}\frac{1}{f(x)}\dd x\le\frac{4}{3}.
\]

\item[四、]（本题满分 12 分）计算三重积分 $\displaystyle\iiint_{(V)}\left(x^{2}+y^{2}\right)\dd V$，其中 $(V)$ 是由
\[
x^{2}+y^{2}+(z-2)^{2}\ge 4,\quad x^{2}+y^{2}+(z-1)^{2}\le 9
\]
及 $z\ge 0$ 所围成的空间图形．

\item[五、]（本题满分 14 分）设 $f(x,y)$ 在区域 $D$ 内可微，且 $\sqrt{\left(\dfrac{\partial f}{\partial x}\right)^{2}+\left(\dfrac{\partial f}{\partial y}\right)^{2}}\le M$，$A(x_{1},y_{1})$，$B(x_{2},y_{2})$ 是 $D$ 内两点，线段 $AB$ 包含在 $D$ 内．证明：
\[
\left|f(x_{1},y_{1})-f(x_{2},y_{2})\right|\le M\,|AB|,
\]
其中 $|AB|$ 表示线段 $AB$ 的长度．

\item[六、]（本题满分 14 分）证明：对于连续函数 $f(x)>0$，有
\[
\ln\int_{0}^{1}f(x)\dd x\ge\int_{0}^{1}\ln f(x)\dd x.
\]

\item[七、]（本题满分 14 分）已知 $\{a_{k}\}$，$\{b_{k}\}$ 是正数数列，且 $b_{k+1}-b_{k}\ge\delta>0$，$k=1,2,\cdots$，$\delta$ 为一常数．证明：若级数 $\displaystyle\sum_{k=1}^{+\infty}a_{k}$ 收敛，则级数
\[
\sum_{k=1}^{+\infty}\frac{k\sqrt[k]{\left(a_{1}a_{2}\cdots a_{k}\right)\left(b_{1}b_{2}\cdots b_{k}\right)}}{b_{k+1}b_{k}}
\]
收敛．
\end{problems}
